Milch und Kaffee haben etwa die Dichte $\rho$ und die spezifische Wärme $c$ von Wasser; Wärmeverluste an Tasse und Luft sind vernachlässigt. Der Kaffee kühlt von $\vartheta_\text{K}$ auf $\vartheta$ ab, die Milch erwärmt sich von $\vartheta_\text{M}$ auf $\vartheta$:
\begin{align}
 c\,\rho V_\text{K}\left(\vartheta_\text{K}-\vartheta\right) &= c\,\rho V_\text{M}\left(\vartheta-\vartheta_\text{M}\right)
\end{align}
Aufgelöst nach dem Milchvolumen:
\begin{align}
 V_\text{M} &= \VMF \\
   &= \VKO\cdot\frac{\TKO-\TmO}{\TmO-\TMO} \\
   &= \VMl \approx \result{\VMlP{-3}{2}}
\end{align}
